% IAC 2026 manuscript stub — claim/evidence plan lives in PLAN.md + CLAIM_EVIDENCE_LEDGER.md.
% Do not rewrite Results numbers here until ledger support levels are closed.
% Official IAF: Word template only; this uses article + iac2026.sty.
\documentclass[10pt,twocolumn,letterpaper]{article}
\usepackage{iac2026}

\renewcommand{\IACpapercode}{IAC-26,A3,IP,xxx,xXXXXX}

\begin{document}

\IACmaketitle
{ARTPS: Operationally Safe Target Prioritization via Depth-Enhanced Hybrid Anomaly Detection}
{Poyraz Baydemir}
{Sel\c{c}uk University, Konya, T\"urkiye}
{%
Planetary rovers operate under strict bandwidth and energy constraints with delayed ground
feedback, which motivates reliable onboard prioritization of scientific targets. Visual-only
autonomy is brittle in field imagery: illumination changes, cast shadows, and low-texture
regions can trigger false positives, while diversity heuristics can suppress high-value
anomalies that appear similar to previously selected targets.
This paper presents ARTPS, an edge-oriented framework that couples monocular relative-depth
estimation with multi-cue anomaly detection and robust localization.
ARTPS fuses image and depth cues via entropy-weighted dynamic fusion, ranks candidates with
a curiosity score and soft similarity penalty, and uses a Priority Buffer for second-chance
re-evaluation of high-anomaly targets down-weighted by diversity policy.
Operational gates address rover-body and boundary-shadow false positives, protect objects
inside shadows, and apply an image-relative size--distance policy based on apparent size
(not metric distance).
On NASA Mars rover imagery, ARTPS reports 0.894 AUROC, 0.847 AUPRC, and 0.823 F1 versus a
PaDiM/PatchCore baseline (0.856 AUROC).
In a lightweight configuration excluding learned depth and autoencoder inference, the core
pipeline sustains 28.1~FPS at $256{\times}256$ on a development workstation.
Proxy ablations (not human bounding-box ground truth) indicate reduced shadow-dense false
positives and lower field-scale false-positive rate under the size--distance policy.
}%
{Mars rover; autonomous exploration; anomaly detection; relative depth; target prioritization; operational safety}

\section{Introduction}
% Outline only — full prose in a later PR. See PLAN.md and CLAIM_EVIDENCE_LEDGER.md.
% (i) Problem: bandwidth/energy + delayed ground feedback; brittle visual-only prioritization.
% (ii) Gap: anomaly baselines lack operational FP control, diversity-aware ranking, relative-depth gates.
% (iii) Contributions: entropy fusion; curiosity + soft similarity + Priority Buffer;
%      rover/shadow FP gates + size--distance policy; edge profiles (workstation timings).
% Related Work: fold into this section (short); no separate long RW chapter.
Planetary surface operations motivate onboard screening of scientifically interesting targets
under communications delay~\citep{estlin2012}.
Standard anomaly detectors provide useful baselines but do not by themselves encode
operational false-positive control or relative-depth guidance for field
imagery~\citep{defard2020padim,roth2022patchcore,ranftl2021dpt}.
% TODO: compress Full_Baydemir_ARTPS.pdf §1--§2 without ViT/FFN textbook derivations.

\section{Material and Methods}
% Stub — keep original equations in a later prose PR (see PLAN.md keep-list).
% Architecture: RGB → relative depth → enhance → cues → entropy fusion → localize/merge
% → curiosity + diversity → Priority Buffer.

\subsection{System architecture}
ARTPS processes an RGB frame through relative-depth estimation, optional enhancement,
multi-cue anomaly extraction, entropy-weighted fusion, localization with tightened merge,
and ranked selection with a Priority Buffer.

\subsection{Depth-guided enhancement}
% KEEP equation: depth-guided enhancement.
% Real-ESRGAN: one paragraph only; no Results claim (unsupported without ablation).
Depth-guided enhancement modulates local contrast using the relative-depth map.
Optional Real-ESRGAN may appear as a profile-aware enhance path; it is not an ablation claim.

\subsection{Anomaly cues and entropy-weighted fusion}
% KEEP: entropy-weighted fusion.
Image, depth, and optional learned anomaly maps are combined with entropy-weighted fusion
that adapts component weights at runtime~\citep{defard2020padim,roth2022patchcore}.

\subsection{Localization, false-positive gates, and size--distance policy}
% KEEP: size--distance policy; tight merge; rover/shadow FP; object-in-shadow gate.
Localization uses tightened box merge.
Rover-body and boundary-shadow masks suppress common false positives; an object-in-shadow
gate protects shadowed rock-like support.
Size--distance policy uses image-relative relative-far and apparent size (not metric size).

\subsection{Curiosity score, diversity, and Priority Buffer}
% KEEP: curiosity score; soft similarity; uncertainty/disagreement; Priority Buffer rule.
Candidates are ranked by a curiosity score with a soft similarity penalty.
The Priority Buffer retains high raw-anomaly candidates down-weighted by diversity policy
for second-chance review when resources allow.

\section{Experimental Protocol}
\label{sec:protocol}
This section distinguishes two evaluation tracks that must not be conflated.
The historical track restates accepted-abstract detection and speed numbers.
The current track is a frozen, version-controlled protocol on the present implementation,
with a supplementary human-reviewed label audit.
A qualitative processing example is included only as an illustration of the frozen
image-score path; it is not a third quantitative track.

\subsection{Historical accepted-abstract track}
The accepted IAC abstract reports AUROC~$0.894$, AUPRC~$0.847$, F1~$0.823$, baseline
AUROC~$0.856$, and $28.1$\,FPS at $256\times256$ for a lightweight core.
The performance values reported in the accepted IAC abstract are retained as the historical
ARTPS evaluation.
That track is a prior experimental snapshot; it is not re-run or retuned in this manuscript,
and the present repository snapshot does not independently reconstruct it.
The supplementary evaluation introduced below follows a separate contemporary protocol and
is therefore reported independently rather than treated as a direct reproduction of the
historical experiment.
No new experiment in this manuscript invents replacement headline numbers.

\subsection{Current reproducible evaluation (\texttt{independent\_eval\_v1})}
A separate frozen evaluation protocol, \texttt{independent\_eval\_v1}, was established for
the current implementation.
It is \emph{not} a reproduction of the accepted-abstract numbers.
Protocol definitions, annotation rules, and frozen execution settings are version-controlled
with the implementation, including the image-score aggregation, preprocessing profile, and
checkpoint hashes used for scoring.
The primary image score is the frozen image-level score defined in Methods
(\texttt{max\_valid\_candidate\_after\_masks}).
Curiosity ranking, diversity penalty, and the Priority Buffer remain disabled for this score.
The held-out test partition was not opened in the present work: no test inference, no test
metric, and no threshold retune on held-out images.
Validation artifacts exist; numeric supplementary results below use the human-reviewed
annotation version only.
No fabricated $N$ and no invented test metrics are reported.
The qualitative example in Figure~\ref{fig:qualitative} uses one deterministically selected
non-test image from the same manifest under the frozen image-score path; it is not a
benchmark and does not open the test split.

\subsection{Supplementary human-reviewed annotation (\texttt{independent\_eval\_v1\_1})}
A separate annotation version, \texttt{independent\_eval\_v1\_1}, records a model-blind
repeat-author review of all 360 included images under the same annotation guide.
The original \texttt{independent\_eval\_v1} manifest remains immutable.
Split assignments are unchanged; the test split remains unopened (human viewing only;
no ARTPS inference).
Scores were remapped onto the reviewed labels with no tuning and no inference rerun.
This supplementary track is a current reproducible validation audit.
It is \emph{not} a reproduction of the accepted-abstract numbers, not a replacement for
the historical experiment, and not an independent second annotation.

\subsection{Annotation and dataset plan}
Annotation follows the independent-evaluation annotation guide: operational target labels,
mask validity, and refusal to proceed when label ratios or frozen splits are incomplete.
Dataset construction uses declared train/val/test planning language; this draft does not
state unverified image counts.

\subsection{Baselines}
Planned baselines include PaDiM- and PatchCore-style anomaly
detectors~\citep{defard2020padim,roth2022patchcore}, evaluated separately from ARTPS fusion
and operational modules.
Baseline comparison tables appear only after \texttt{independent\_eval\_v1} metrics exist.

\subsection{Metrics and statistics}
Named metrics for the current protocol include AUROC, AUPRC, F1, and throughput (FPS) where
profiles request them.
For the supplementary \texttt{independent\_eval\_v1\_1} validation, AUROC is the principal
reported metric and threshold-dependent F1 is not treated as a performance measure.
Bootstrap confidence intervals are \emph{planned} once frozen metrics exist; this draft does
not invent interval values.
% TODO(results): fill bootstrap CIs after independent\_eval\_v1 run artifacts.

\subsection{Reproducibility and evaluation integrity}
Evaluation is aborted when required inputs or frozen protocol conditions are incomplete.
Version-controlled protocol definitions alone are not a substitute for a completed run.

\subsection{Hardware profiles}
Workstation timing profiles and a Jetson / edge-hardware protocol are defined at the planning
level.
No Jetson numbers are reported in this draft.
Hardware measurements are not flight hardware and do not constitute flight certification.

\subsection{Proxy analyses (non-headline)}
Preliminary proxy analyses for shadow/FP gates and size--distance policy use curated labels,
masks, and OFF-run pseudo-ground-truth.
They are not human bounding-box ground truth and are excluded from abstract headlines and
from primary Results tables in this draft.

\section{Results}
% NO rewritten result tables in this PR.
% C05--C07: accepted_abstract_reproduction_pending (keep abstract numbers; no final tables).
% C09--C11: preliminary proxy analysis only (not abstract headlines).
% Bind claims to CLAIM_EVIDENCE_LEDGER.md and FIGURE_TABLE_PLAN.md.
Results will summarize detection, ranking, localization, field-condition, and hardware
profiles once ledger rows close beyond accepted-abstract reproduction pending.
Preliminary proxy analyses (shadow/FP and size--distance on a curated non--human-GT set)
may be noted as such; they are not restated as final manuscript headline results here.

\section{Discussion}
\label{sec:discussion}
ARTPS is best read as an operational architecture: anomaly cues and relative depth feed
fusion, localization, suppression, diversity, and buffering under fail-closed gates.
The design goal is structured prioritization under rover constraints, not unconstrained
claim of superior detection in all regimes.

Relative depth is an ordering cue~\citep{ranftl2021dpt}; interpreting it as metric distance
would overstate what monocular estimation provides.
Operational suppressions address common imagery distractors, but residual false positives
remain an open operational risk on the still-closed independent test split.

% TODO(results): discuss accepted-abstract reproduction only if C05--C07 close beyond pending.
% Avoid ``outperforms'' / ``demonstrate'' language until evidence gates allow.

The separation between historical accepted-abstract claims and the current independent-eval
track is intentional: conflating them would misrepresent evidence status.
The initial \texttt{independent\_eval\_v1} binary labels were produced by a classical-image
heuristic with forced $180/180$ class balance.
A model-blind repeat-author review changed the semantic distribution substantially
($340$ positive / $20$ negative).
\texttt{independent\_eval\_v1\_1} therefore follows the intended science-interest annotation
definition more faithfully on this selected domain; the imbalance indicates that
science-interest prevalence is naturally high under that definition.
Frozen-score AUROC~$0.772$ indicates useful ranking separation.
Threshold-based F1 is not informative here because the predeclared operating point is
all-positive.
A future balanced challenge set or harder negative set may be useful for operating-point
evaluation.
This label-audit finding is confined to the supplementary independent-eval benchmark; it
does not imply that the historical accepted-abstract ARTPS experiment suffered the same
issue.
Proxy analyses remain secondary and must not migrate into abstract headlines.

Safety language stays bounded: ARTPS is safety-aware software on a flight-readiness-oriented
path; it is not flight-qualified and does not assert universal safety guarantees.
Hardware timing (including planned Jetson profiles) is engineering measurement, not
certification evidence.

\section{Conclusions}
\label{sec:conclusion}
ARTPS specifies an onboard target-prioritization architecture that couples multi-cue visual
anomaly responses with monocular relative depth, entropy-weighted fusion, operational
suppression, apparent size--distance policy, diversity ranking, and a Priority Buffer.

Accepted-abstract detection and speed numbers are retained as historical claims pending
repository-level reproduction.
A separate current evaluation protocol (\texttt{independent\_eval\_v1}) remains test-closed;
a supplementary human-reviewed validation audit (\texttt{independent\_eval\_v1\_1}) reports
ranking-oriented AUROC~$0.772$ under frozen scores and does not open the test split.

Future work includes a harder-negative or independently annotated challenge set, optional
reproduction of historical abstract claims if artifacts allow, planned edge-hardware timing
profiles, and mission-specific V\&V beyond the scope of this software manuscript.


\bibliographystyle{unsrtnat}
\bibliography{references}

\end{document}
